\documentclass[10pt,a4paper]{article}
\usepackage[utf8]{inputenc}
\usepackage[vietnamese]{babel}
\usepackage{amsmath,amssymb,amsfonts}
\usepackage{geometry}
\geometry{a4paper, left=10mm, right=10mm, top=10mm, bottom=10mm}
\usepackage{tikz}
\usetikzlibrary{shapes,arrows.meta,positioning,calc}
\usepackage[most]{tcolorbox}
\usepackage{array,booktabs,multicol}
\usepackage{fontawesome}
\usepackage{helvet}
\renewcommand{\familydefault}{\sfdefault}
\usepackage{parskip}

\definecolor{primary}{RGB}{112, 48, 160}
\definecolor{teal}{RGB}{0, 128, 128}
\definecolor{orange}{RGB}{230, 108, 44}
\definecolor{green}{RGB}{46, 139, 87}
\definecolor{rose}{RGB}{192, 0, 0}

\pagestyle{empty}

\begin{document}

% --- HEADER BANNER ---
\begin{tcolorbox}[
  enhanced,
  colback=violet!8!white,
  colframe=violet!80!black,
  arc=2.5mm, boxrule=1pt,
  left=3mm, right=3mm, top=2mm, bottom=2mm
]
  \begin{tabular*}{\textwidth}{@{\extracolsep{\fill}} lr}
    {\small\bfseries\color{violet!85!black}\faBook\ SỔ TAY KIẾN THỨC TOÁN 2} & 
    {\small\bfseries\color{teal!85!black}\faGraduationCap\ CHƯƠNG TRÌNH GDPT 2018}
  \end{tabular*}
  \vspace{-1mm}
  \begin{center}
    \rule{0.98\textwidth}{0.6pt}
  \end{center}
  \vspace{-2mm}
  \centering
  {\LARGE\bfseries\color{violet!90!black}\faBookmark\ BÀI 60: PHÉP CỘNG (CÓ NHỚ) TRONG PHẠM VI 1000}\\[1mm]
  {\itshape\footnotesize\color{black!75} Kĩ thuật cộng có nhớ sang hàng chục hoặc sang hàng trăm}
\end{tcolorbox}

\vspace{-2mm}

\noindent
\begin{minipage}[t]{0.49\textwidth}\vspace{0pt}
\begin{tcolorbox}[
  enhanced,
  colback=violet!4!white,
  colframe=violet!75!black,
  arc=2mm, boxrule=0.8pt,
  left=2mm, right=2mm, top=1.8mm, bottom=1.8mm,
  title={\footnotesize\bfseries\color{white}\faKey\ 1. Quy tắc cộng có nhớ}
]
{\footnotesize
\begin{itemize}\setlength{\itemsep}{1pt}\setlength{\leftskip}{-2.5mm}
  \item Khi cộng hai chữ số ở một hàng mà kết quả $\ge 10$:
  \begin{itemize}\setlength{\itemsep}{0.5pt}\setlength{\leftskip}{-2mm}
    \item Ta viết chữ số hàng đơn vị của kết quả.
    \item \textbf{Nhớ $1$} sang hàng tiếp theo ở bên trái.
  \end{itemize}
  \item \textbf{Quy tắc tính toán:}
  \begin{itemize}\setlength{\itemsep}{0.5pt}\setlength{\leftskip}{-2mm}
    \item Nếu nhớ từ hàng đơn vị sang hàng chục: cộng thêm $1$ vào tổng các chữ số hàng chục.
    \item Nếu nhớ từ hàng chục sang hàng trăm: cộng thêm $1$ vào tổng các chữ số hàng trăm.
  \end{itemize}
  \item Trong chương trình lớp 2, ta học phép cộng có nhớ \textbf{không quá 1 lần}.
\end{itemize}
}
\end{tcolorbox}
\end{minipage}
\hfill
\begin{minipage}[t]{0.49\textwidth}\vspace{0pt}
\begin{tcolorbox}[
  enhanced,
  colback=teal!4!white,
  colframe=teal!75!black,
  arc=2mm, boxrule=0.8pt,
  left=2mm, right=2mm, top=1.8mm, bottom=1.8mm,
  title={\footnotesize\bfseries\color{white}\faSliders\ 2. Hai trường hợp có nhớ điển hình}
]
{\footnotesize
\begin{itemize}\setlength{\itemsep}{1pt}\setlength{\leftskip}{-2.5mm}
  \item \textbf{Trường hợp 1: Nhớ sang hàng chục}\\
  Ví dụ: $247 + 325$:
  \begin{itemize}\setlength{\itemsep}{0.5pt}\setlength{\leftskip}{-2mm}
    \item $7 + 5 = 12$, viết $2$ \textbf{nhớ $1$}.
    \item $4 + 2 = 6$, thêm $1$ bằng $\mathbf{7}$, viết $7$.
    \item $2 + 3 = 5$, viết $5$.
    \item Kết quả: $\mathbf{247 + 325 = 572}$.
  \end{itemize}
  \item \textbf{Trường hợp 2: Nhớ sang hàng trăm}\\
  Ví dụ: $382 + 254$:
  \begin{itemize}\setlength{\itemsep}{0.5pt}\setlength{\leftskip}{-2mm}
    \item $2 + 4 = 6$, viết $6$.
    \item $8 + 5 = 13$, viết $3$ \textbf{nhớ $1$}.
    \item $3 + 2 = 5$, thêm $1$ bằng $\mathbf{6}$, viết $6$.
    \item Kết quả: $\mathbf{382 + 254 = 636}$.
  \end{itemize}
\end{itemize}
}
\end{tcolorbox}
\end{minipage}

\vspace{-1.5mm}
\begin{tcolorbox}[
  enhanced,
  colback=orange!4!white,
  colframe=orange!80!black,
  arc=2mm, boxrule=0.8pt,
  left=2mm, right=2mm, top=1.5mm, bottom=1.5mm,
  title={\footnotesize\bfseries\color{white}\faSitemap\ 3. Sơ đồ xử lí nhớ sang hàng liền kề bên trái}
]
\centering
\begin{tikzpicture}[x=1.65cm, y=0.75cm]
  \tikzset{arr/.style={-{Stealth[scale=0.8]}}, every node/.append style={align=center}}
  \node[rectangle, rounded corners=2pt, draw=violet!80!black, fill=violet!15, font=\scriptsize\bfseries] (dv) at (0,0) {HÀNG ĐƠN VỊ\\$7 + 5 = 12$};
  \node[rectangle, rounded corners=2pt, draw=teal!80!black, fill=teal!15, font=\scriptsize\bfseries] (ch) at (2.4,0) {HÀNG CHỤC\\$4 + 2 = 6$};
  \node[rectangle, rounded corners=2pt, draw=green!70!black, fill=green!15, font=\scriptsize\bfseries] (tr) at (4.8,0) {HÀNG TRĂM\\$2 + 3 = 5$};

  \draw[arr, very thick, red!85!black] (dv) -- node[above, font=\tiny\bfseries] {Nhớ $1$} node[below, font=\tiny] {Viết $2$} (ch);
  \node[font=\tiny\bfseries, color=teal!90!black] at (2.4,-0.6) {Thêm $1 \implies 6 + 1 = 7$};
\end{tikzpicture}
\end{tcolorbox}

\vspace{-1.5mm}
\noindent
\begin{minipage}[t]{0.49\textwidth}\vspace{0pt}
\begin{tcolorbox}[
  enhanced,
  colback=green!4!white,
  colframe=green!60!black,
  arc=2mm, boxrule=0.8pt,
  left=2mm, right=2mm, top=1.8mm, bottom=1.8mm,
  title={\footnotesize\bfseries\color{white}\faCheckCircle\ 4. Ví dụ thực hành đặt tính}
]
{\footnotesize
\textbf{Đặt tính rồi tính:}\\
\begin{tabular}{r@{\,}c@{\,}lcr@{\,}c@{\,}l}
  & $436$ & & & & $584$ & \\
$+$ & $218$ & & &$+$ & $142$ & \\
\cline{1-3}\cline{5-7}
  & $\mathbf{654}$ & & & & $\mathbf{726}$ &
\end{tabular}
\vspace{1mm}\\
- $6 + 8 = 14$, viết $4$ nhớ $1$; $3 + 1 = 4$, thêm $1$ là $\mathbf{5}$; $4 + 2 = \mathbf{6}$.\\
- $4 + 2 = 6$, viết $6$; $8 + 4 = 12$, viết $2$ nhớ $1$; $5 + 1 = 6$, thêm $1$ là $\mathbf{7}$.
}
\end{tcolorbox}
\end{minipage}
\hfill
\begin{minipage}[t]{0.49\textwidth}\vspace{0pt}
\begin{tcolorbox}[
  enhanced,
  colback=rose!4!white,
  colframe=rose!70!black,
  arc=2mm, boxrule=0.8pt,
  left=2mm, right=2mm, top=1.8mm, bottom=1.8mm,
  title={\footnotesize\bfseries\color{white}\faExclamationTriangle\ 5. Lỗi quên cộng số nhớ}
]
{\footnotesize
\begin{itemize}\setlength{\itemsep}{1.5pt}\setlength{\leftskip}{-2.5mm}
  \item \textbf{Sai sót phổ biến nhất:} Nhớ ở hàng đơn vị nhưng sang hàng chục lại \textbf{quên không cộng thêm 1}.
  \item \textbf{Mẹo khắc phục:} Có thể chấm một dấu chấm nhỏ ở phía trên chữ số hàng chục để nhắc nhở bản thân nhớ cộng thêm $1$.
  \item Luôn nhẩm lại: \textit{"nhớ 1 thì phải cộng 1"}.
\end{itemize}
}
\end{tcolorbox}
\end{minipage}


\end{document}
